\chapter{Verificación y Evaluación Experimental}
\label{ch:verificacion}

Este capítulo describe la estrategia de verificación y las actividades de
evaluación experimental realizadas para el \textbf{Subsystema de
Super--Resolución basada en Inteligencia Artificial (SAI)}. La verificación se
enfoca en comprobar el correcto funcionamiento del pipeline de procesamiento, la
consistencia de los resultados generados y el cumplimiento de los requerimientos
definidos en el Capítulo~\ref{ch:requerimientos}.

La evaluación se plantea desde un enfoque cualitativo y comparativo, dado que el
subsistema SAI se concibe como una herramienta de apoyo visual y no como un
instrumento de medición científica primaria.

% -------------------------------------------------------------------
\section{Estrategia general de verificación}

La verificación del SAI se basa en los siguientes principios:

\begin{itemize}
    \item Uso de imágenes reales adquiridas por el CMOS Microscopy Payload.
    \item Comparación directa entre imagen original y salida procesada.
    \item Preservación de trazabilidad entre datos de entrada y salida.
    \item Evaluación visual y métrica básica sin interpretación metrológica.
\end{itemize}

No se emplean métricas de resolución absoluta, dado que el procesamiento no
incrementa la resolución óptica real del sistema.

% -------------------------------------------------------------------
\section{Conjunto de datos de prueba}

Las pruebas se realizaron utilizando imágenes experimentales capturadas en
entorno de laboratorio, correspondientes a muestras biológicas planas
(ej. epidermis de cebolla) adquiridas en configuración de microscopía sin lentes.

Las imágenes utilizadas presentan:
\begin{itemize}
    \item Resolución nativa del sensor CMOS.
    \item Iluminación uniforme controlada.
    \item Parámetros de captura registrados como metadatos.
\end{itemize}

Este conjunto de datos representa escenarios reales de operación del sistema.

% -------------------------------------------------------------------
\section{Procedimiento de verificación funcional}

El procedimiento general de verificación se describe a continuación:

\begin{enumerate}
    \item Selección de imagen original adquirida por el CMOSM.
    \item Carga de la imagen en el pipeline SAI.
    \item Ejecución del preprocesamiento definido.
    \item Aplicación del modelo de super--resolución seleccionado.
    \item Generación de la imagen procesada.
    \item Almacenamiento separado de la imagen original y la imagen procesada.
    \item Registro de metadatos del procesamiento.
\end{enumerate}

Cada ejecución se documentó para asegurar repetibilidad del proceso.

% -------------------------------------------------------------------
\section{Criterios de evaluación}

La evaluación de resultados se realizó considerando los siguientes criterios:

\subsection{Consistencia estructural}
Se verificó que las estructuras visibles en la imagen original se mantuvieran
presentes en la imagen procesada, sin desaparición o deformación significativa
de bordes.

\subsection{Ausencia de artefactos críticos}
Se inspeccionó visualmente la presencia de:
\begin{itemize}
    \item patrones artificiales no presentes en la imagen original,
    \item halos excesivos alrededor de bordes,
    \item texturas irreales inducidas por el modelo.
\end{itemize}

\subsection{Mejora perceptual}
Se evaluó la mejora subjetiva en términos de:
\begin{itemize}
    \item definición de bordes,
    \item legibilidad de estructuras celulares,
    \item contraste local.
\end{itemize}

Esta evaluación se realizó de manera comparativa y cualitativa.

% -------------------------------------------------------------------
\section{Resultados experimentales}

Los resultados obtenidos muestran que el subsistema SAI:

\begin{itemize}
    \item mejora la legibilidad visual de estructuras ya presentes en la imagen
    original,
    \item no introduce artefactos dominantes cuando se emplean parámetros
    conservadores,
    \item preserva la correspondencia espacial con la imagen de entrada.
\end{itemize}

En particular, se observó una mejora perceptual clara en la visualización de
paredes celulares en muestras vegetales, facilitando la interpretación visual
sin modificar la información base adquirida.

La Figura~\ref{fig:resultados_esrgan} presenta ejemplos representativos de los
resultados experimentales obtenidos mediante la aplicación del modelo ESRGAN x4
sobre imágenes adquiridas por el CMOSM.

\begin{figure}[H]
    \centering
    \begin{subfigure}[b]{0.48\textwidth}
        \centering
        \includegraphics[width=\textwidth]{figures/Comparación SSR 2.png}
        \caption{Comparación 1: Super-resolución x4 aplicada a tejido vegetal}
        \label{fig:esrgan_test1}
    \end{subfigure}
    \hfill
    \begin{subfigure}[b]{0.48\textwidth}
        \centering
        \includegraphics[width=\textwidth]{figures/Comparación SSR.png}
        \caption{Comparación 2: Super-resolución x4 mostrando mejora en definición}
        \label{fig:esrgan_test2}
    \end{subfigure}
    \caption{Resultados experimentales de super-resolución mediante ESRGAN x4, demostrando mejora en definición de bordes y legibilidad de estructuras celulares sin introducir artefactos dominantes.}
    \label{fig:resultados_esrgan}
\end{figure}

% -------------------------------------------------------------------
\section{Verificación de requerimientos}

La Tabla~\ref{tab:verificacion-sai} resume el estado de verificación de los
requerimientos principales del subsistema SAI.

\begin{table}[H]
\centering
\begin{tabular}{p{3cm} p{7cm} p{3cm}}
\toprule
\textbf{ID} & \textbf{Requerimiento} & \textbf{Estado} \\
\midrule
FR-SAI-01 & Ingesta de imágenes del CMOSM & Verificado \\
FR-SAI-02 & Aplicación de super--resolución & Verificado \\
FR-SAI-03 & Preservación de imagen original & Verificado \\
PR-SAI-04 & Ausencia de artefactos críticos & Parcialmente verificado \\
IR-SAI-01 & Compatibilidad de formatos & Verificado \\
\bottomrule
\end{tabular}
\caption{Estado de verificación de requerimientos del SAI.}
\label{tab:verificacion-sai}
\end{table}

Los requerimientos marcados como parcialmente verificados corresponden a
criterios perceptuales que requieren mayor número de muestras para validación
estadística.

% -------------------------------------------------------------------
\section{Limitaciones de la evaluación}

La evaluación experimental presenta las siguientes limitaciones:

\begin{itemize}
    \item Tamaño reducido del conjunto de datos.
    \item Dependencia del modelo entrenado y sus parámetros.
    \item Evaluación principalmente cualitativa.
\end{itemize}

Estas limitaciones se consideran aceptables dada la naturaleza exploratoria del
subsistema SAI.

% -------------------------------------------------------------------
\section{Resumen}

La verificación realizada demuestra que el subsistema SAI cumple su función como
herramienta de apoyo visual para la interpretación de imágenes adquiridas por el
CMOSM. Los resultados obtenidos justifican su integración como módulo
complementario dentro del proyecto INTISAT, manteniendo un enfoque conservador
respecto a su uso científico.
